{\large\textbf{QUESTION T2 [10 pts]. Cooking with sunlight?}}

The redistribution of luminous flux into regions of high and low intensity due
to reflection or refraction can lead to the formation of \textit{caustics:}
curves formed by the intersection of infinitesimally close light rays. As an
example, Figure 2a shows the caustic formed by rays incident paraxially on a
half-cylindrical mirror. Figure 2b shows a photograph of the caustic formed by
the reflection of light on the bottom part of a coffee cup. The brightest and
central point of the caustic is called the \textit{cusp}; at this point, the
light intensity is maximal.

\begin{center}
\begin{minipage}[b]{0.49\linewidth}\centering
\includegraphics[width=\linewidth]{figures/IPhO_2026_Q2_fig1.png}

\texttt{Fig. 2a.}
\end{minipage}\hfill
\begin{minipage}[b]{0.40\linewidth}\centering
\includegraphics[width=0.80\linewidth]{figures/IPhO_2026_Q2_fig2.jpg}

\texttt{Fig. 2b.}
\end{minipage}
\end{center}

Caustics have a variety of applications in science, ranging from
high-precision optics and quantum optics to astrophysics. In this question, you
will study the caustic formation by light reflection on a half-cylinder mirror
of radius $R$, as shown in figure 2c. You will also investigate other aspects of
this mirror, such as its application in \textit{solar cookers}.

Throughout the problem, assume that the incident rays come from a distant
external light source and are parallel to the mirror's optical axis. By
symmetry, it is sufficient to study reflection in one plane. Figure 2d shows
such a plane, the mirror, and the coordinate system used throughout the
problem.

\begin{center}
\begin{minipage}[b]{0.36\linewidth}\centering
\includegraphics[width=0.78\linewidth]{figures/IPhO_2026_Q2_fig3.png}

\texttt{Fig. 2c.}
\end{minipage}\hfill
\begin{minipage}[b]{0.56\linewidth}\centering
\includegraphics[width=\linewidth]{figures/IPhO_2026_Q2_fig4.png}

\texttt{Fig. 2d.}
\end{minipage}
\end{center}

\textbf{T2-A [1.5 pts]. Multiple reflections}

\begin{center}\includegraphics[width=0.48\linewidth]{figures/IPhO_2026_Q2_fig5.png}\end{center}

\begin{center}\texttt{Fig.2e.}\end{center}

Figure 2e shows the relationship between the $x$-coordinate of an incident ray
and the number $N$ of reflections it undergoes, up to $N = 3$. As $x$ varies
over the interval $(-R, R)$, the variable $N$ varies over all natural numbers
greater than or equal to 1. The maximum distance from the optical axis that
allows an incident ray to be reflected by the mirror at most $N$ times is
denoted $x_N$.

\textbf{T2-A1 [1.5 pts].} Find the general expression for $x_N$ in the
infinite sequence $x_1, x_2, x_3, \ldots$ that limits the coordinates of the
rays with at most $N$ reflections. Express your answer in terms of $R$ and $N$.

\textbf{T2-B [4.0 pts]. Solar Cooker}

\begin{center}\includegraphics[width=0.50\linewidth]{figures/IPhO_2026_Q2_fig6.png}\end{center}

\begin{center}\texttt{Fig.2f.}\end{center}

Consider a system consisting of a half-hollow-cylinder mirror and a cylindrical
container. The cylinder is fully absorbing: any ray of light hitting it will be
absorbed. The mirror radius is $R$. The container radius $a$. The axes of the
mirror and the container are parallel. The center of the container is $R/2$
from the mirror on the system's symmetry plane, as shown in figure 2f. The
mirror harnesses solar radiation to heat food inside the container. Assume
that the sunlight is of constant and uniform intensity (power per unit area),
and that the light rays are parallel to the optical axis of the mirror.
Furthermore, assume that $a$ is such that any ray that is absorbed by the
cylinder reflects from the mirror at most once. Let $\theta_{\mathrm{max}}$ be
the maximum angle of incidence on the mirror (measured with respect to the
normal drawn at the point of incidence) of any reflected ray striking the
container.

\textbf{T2-B1 [2.0 pts].} The container radius is
$a = \alpha\sin(\theta_{max}) + \beta\sin(2\theta_{max})$. Write $\alpha$ and
$\beta$ in terms of $R$.

If you do not find $\alpha$ and $\beta$, use the expression given for the rest
of the problem.

\textbf{T2-B2 [1.5 pts].} Let $P_0$ be the power that would be received by the
cyllinder if the mirror was not present. Write $P/P_0$, the ratio of the actual
received power $P$ to $P_0$, in terms of $\theta_{\mathrm{max}}$.

\textbf{T2-B3 [0.5 pts].} If $R = 1.0\,\mathrm{m}$, write the value of $a$ such
that $P = 5P_0$. Give your answer in cm.

\textbf{T2-C [4.5 pts]. Caustics and Cusp}

\textbf{T2-C1 [0.5 pts].} Consider a ray $A$ that strikes the mirror at an
angle $\theta$, as shown in Figure 2g. Upon reflection, the equation
corresponding to this ray in the coordinate system defined in Figure 2g. is
$y = m_Ax + b_A$. Write $m_A$ and $b_A$ in terms of $\theta$ and $R$

Hint: the expected form of $m_A = K_1\cot(K_2\theta)$ and
$b_A = R\frac{K_3}{\cos(K_4\theta)}$, with $K_i$ numerical constants.

\begin{center}\includegraphics[width=0.49\linewidth]{figures/IPhO_2026_Q2_fig7.png}\end{center}

\begin{center}Fig. 2g.\end{center}

\textbf{T2-C2 [2.0 pts].} Now consider a ray $B$, parallel to $A$, that strikes
the surface at an angle $\theta + \Delta\theta$, with
$\Delta\theta \ll \theta$. Upon reflection, the equation corresponding to this
ray is $y = m_Bx + b_B$. Write $m_B$ and $b_B$ in terms of $\theta$,
$\Delta\theta$, and $R$, approximating to first order in $\Delta\theta$.

\textbf{T2-C3 [1.0 pts].} Find the coordinates of the point of intersection
$(X_c, Y_c)$ between the reflected rays of $A$ and $B$. Write your answer in
terms of $R$ and $\theta$.

\textbf{T2-C4 [1.0 pts].} For $\theta \ll 1$, we have
$Y_c = v|X_c|^{p/q}+u$. Write $v$ and $u$ in terms of $R$ and find the values
of the integers $p$ and $q$ (equivalently, give the rational number $p/q$).
